\section{读懂 Scaling Law：Fit、Residual 与 Compute Allocation}

前一节解释了研究路径，本节进入公式与决策。常见语言模型 scaling law 把 cross-entropy loss 写成资源的幂律衰减加不可约项。重点不是背 exponent，而是理解每个 fit 假设、数据范围与误差如何影响数千万乃至更高成本的 compute allocation。

\subsection{Power Law 与不可约损失}

\term{cross-entropy loss}（交叉熵损失）衡量模型给真实下一个 token 分配的负对数概率。一个简化 compute scaling 形式是
\[
L(C)=L_{\infty}+aC^{-\alpha},
\]
其中 $C$ 表示训练 compute，$L(C)$ 是对应 loss，$L_{\infty}$ 是当前数据/任务假设下的不可约项，$a$ 是尺度系数，$\alpha>0$ 是收益衰减速度。Log-log 图可能近似直线，但 $L_{\infty}$、噪声与 recipe change 会显著影响外推。

\lecturefigure{03-scaling-law-fit.png}{Scaling law 是带 residual、confidence 和适用边界的 fit，而不是永恒直线。}{Kaplan et al. 2020；本地字幕 06:20--11:50；概念重绘。}

\begin{knowledgebox}{读图：先看残差，再看斜率}
Axes 必须说明是 FLOPs、token、parameter 还是 wall-clock；fit 要给训练区间和 confidence；residual 显示 architecture、data 或 optimization 是否系统偏离；decision 才能据此调整 budget、model/data mix 或停止外推。若较大 run 连续落在预测线同一侧，最重要的信息可能不是“模型还不错”，而是当前 law 已失配。
\end{knowledgebox}

\begin{warningbox}{Benchmark 提升与 loss scaling 不是同一函数}
Average loss 的平滑下降可以对应某些 downstream task 的突然阈值效应，也可能不改善目标任务。能力评估要覆盖真实环境、工具和 elicitation；不能用一条 pre-training loss 曲线替代 post-training quality、安全和 serving economics。
\end{warningbox}

\subsection{Compute multiplier 是组合，而不是秘密常数}

课堂把 architecture、data、optimizer 和 systems improvement 统称为 compute multiplier：同样 nominal FLOPs 下获得更多 effective capability。保密具有竞争理由，但内部仍要能复现。每个 multiplier 应有 baseline、ablation、interaction 与 transfer evidence，避免 portfolio 中多个改动相互抵消却都被宣称成功。

\lecturefigure{04-compute-allocation.png}{Compute allocation 同时投资 model、data、optimization 与 systems 四类 multiplier。}{本地字幕 12:00--15:20；概念重绘。}

\begin{knowledgebox}{读图：收益要按 effective compute 归因}
Architecture 可能提高表达或稀疏效率；data 改善 signal-to-noise；optimization 提高稳定性与收敛；systems 提高 hardware utilization。最终应比较固定 wall-clock、固定美元、固定 energy 或固定 FLOPs 下的 quality，而不是只报告某一 microbenchmark。Multiplier 还可能只在特定 scale 生效，必须验证 transfer。
\end{knowledgebox}

\teachervoice{讲者把 frontier lab 描述成 researchers 与 engineers 深度协同，而不是研究员提出想法、工程师做“杂活”。Scaling law 让两者共享一个实验语言：研究改 recipe，工程保证 run 可比、可恢复、可测量。}

\subsection{本章小结}

Scaling law 是资源决策工具。正确做法是联合 fit、residual、uncertainty 与 multiplier ablation；错误做法是把历史指数当自然常数，或把系统吞吐提升直接等同于模型能力提升。

